\section{Sparse assignments from nonamenability}\label{sec:sparse}

The obstruction requires uniform operator bounds together with increasing
energy per fibre dimension. This section supplies its combinatorial
input: a list of convolution edges that can be assigned to their
endpoints with small loads. The matrix weights in \Cref{sec:weights}
will convert those loads into operator bounds.

Write \((R_tf)(x)=f(xt)\) for right translation on \(\ell^2(G)\).
Thus \(R_sR_t=R_{st}\). We first obtain the finite averaging operator
needed for the construction. This is the elementary averaging-norm
form of the random-walk criterion for nonamenability; compare
\cite{Kesten1959}. We give the implication at the precise strength
used here.

\begin{lemma}\label{lem:contracting-average}
If \(G\) is a nonamenable discrete group, there are a finite set
\(S\subset G\) containing \(e\), with \(d=|S|\ge2\), and a number
\(0<\rho<1\) such that
\[
 \norm{\frac1d\sum_{t\in S}R_t}\le\rho.
\]
\end{lemma}

\begin{proof}
Every such average is a contraction. Suppose all have norm one.
For a unit vector \(f\) and a finite \(S\) of cardinality \(d\),
\[
 \norm{\frac1d\sum_{t\in S}R_tf}_2^2
 =1-\frac1{2d^2}\sum_{t,t'\in S}\norm{R_tf-R_{t'}f}_2^2.
\]
Taking vectors on which the average has norm arbitrarily close to one,
and using \(e\in S\), gives unit vectors arbitrarily close to invariant
under any prescribed finite set of right translations.

For each such vector define the mean
\[
 m_f(b)=\sum_{x\in G}|f(x)|^2b(x^{-1}),\qquad b\in\ell^\infty(G).
\]
Changing variables \(y=xg\) gives
\(m_f(L_gb)=m_{R_{g^{-1}}f}(b)\). The Cauchy--Schwarz inequality yields
\[
 |m_f(L_gb)-m_f(b)|
 \le2\norm{b}_\infty\norm{R_{g^{-1}}f-f}_2.
\]
Direct the finite translation tests and their positive error tolerances
by refinement. A weak-star cluster point of the resulting means exists
by compactness of the dual unit ball. It is positive, normalized, and
left invariant by the displayed estimate. This contradicts
nonamenability. Some average therefore has norm less than one; choose
\(\rho\) between its norm and one. The set has at least two elements,
since the average for \(\{e\}\) is the identity.
\end{proof}

The next argument converts an operator-norm estimate into an assignment
of edges with bounded loads at their endpoints. It is related to the
rectangular kernel estimates for Littlewood functions
\cite[Proposition~2.6]{GGMT2018}; we include the complete assignment
argument and retain all multiplicities. The conversion from endpoint
assignments to row--column bounds also parallels the forest orientation
in \cite[Propositions~2.3--2.4]{EpsteinMonod2009}.

\begin{proposition}\label{prop:sparse-splitting}
Let \(G\) be countable and nonamenable, and fix \(S,d,\rho\) as in
\Cref{lem:contracting-average}. For every \(\ell\ge1\), list the products
of \(\ell\) elements of \(S\), with multiplicity, as
\(s_1,\ldots,s_n\), where \(n=d^\ell\), and put
\(r=\lceil n\rho^\ell\rceil\). There is a function
\(a:G\times\{1,\ldots,n\}\to\{0,1\}\) such that
\begin{equation}\label{eq:sparse-mask}
 \#\{i:a(x,i)=1\}\le r\quad(x\in G),\qquad
 \#\{i:a(ys_i^{-1},i)=0\}\le r\quad(y\in G).
\end{equation}
\end{proposition}

\begin{proof}
The product list satisfies
\[
 \sum_{i=1}^nR_{s_i}=\left(\sum_{t\in S}R_t\right)^\ell,
 \qquad
 \norm{\sum_{i=1}^nR_{s_i}}\le n\rho^\ell.
\]
Consider a bipartite multigraph with row and column vertices in
disjoint copies of \(G\). Its labelled edge \((x,i)\) joins row \(x\)
to column \(xs_i\). For finite row and column sets \(U,V\), pairing
the convolution operator with their indicator functions gives
\begin{equation}\label{eq:edge-count}
 \#\{(x,i):x\in U,\ xs_i\in V\}
 \le n\rho^\ell\sqrt{|U||V|}.
\end{equation}

Give each vertex \(r\) distinct slots. We match every labelled edge to
one slot at one of its two endpoints. For a finite edge set \(F\), let
\(U,V\) be its endpoint sets. Its neighboring slots number
\(r(|U|+|V|)\), whereas \Cref{eq:edge-count} gives
\[
 |F|\le n\rho^\ell\sqrt{|U||V|}
 \le\frac r2(|U|+|V|)\le r(|U|+|V|).
\]
This applies to every finite edge set and hence verifies Hall's
condition for every finite family of edges and all its subfamilies.
For completeness, the finite form of Hall's marriage theorem
\cite{Hall1935} follows here by augmenting paths. Take a matching of maximum
cardinality for a finite family. If an edge-vertex remains unmatched,
follow alternating paths from it. A reachable unmatched slot allows
an augmentation. Otherwise every reachable slot is matched to a
reachable edge-vertex other than the initial one. All neighbors of the
reachable edge-vertices are reachable slots, contradicting Hall's
inequality. Hence every finite family has an injective slot assignment.

Enumerate the countable edge set. The assignments for its finite
initial segments form a tree under restriction, with a node at every
depth and at most \(2r\) extensions of each node. Starting at the root,
choose at each step a child with descendants at arbitrarily large
depths; one exists because there are finitely many children. The
resulting infinite branch supplies a simultaneous assignment. Set \(a(x,i)=1\) when the edge is assigned to
its row and zero when it is assigned to its column. No vertex receives
more than \(r\) edges, proving \Cref{eq:sparse-mask}. Parallel edges
remain distinct throughout the matching, and the two endpoints remain
distinct vertices even if \(s_i=e\).
\end{proof}
